% Abstract — working-draft version. The empirical sentences will be finalized
% after the registered held-out comparison on the expanded dataset; nothing
% here claims results that do not yet exist.
\begin{abstract}
What pacing strategy minimizes 200-yard freestyle race time, and do
competitive swimmers use it? We formulate the four-lap race as a constrained
optimization problem over split velocities, with a cost function derived from
measured hydrodynamic drag, a two-parameter aerobic--anaerobic energy system,
and oxygen-uptake kinetics, and we prove two structural results. First, under
a fixed energy budget with velocity-only cost, even pacing is the unique
optimum --- a theorem, holding for every cost exponent $p>1$. Second, the
optimal pacing \emph{shape} is invariant to aerobic capacity, anaerobic
capacity, drag, and efficiency: those parameters set how fast the race is
swum, and provably not how it is divided. Because real 200-yard races are not
evenly paced, we formalize five competing fatigue mechanisms (constant
economy, oxygen kinetics, reserve-coupled cost, position-coupled cost, and a
fatigue-coupled velocity ceiling) that make three qualitatively distinct,
falsifiable predictions about split shape, and we test them against
anonymized official race results from male swimmers aged 15--18 under a
pre-registered analysis plan with a swimmer-grouped held-out test set. In the
pilot sample (80 races, one meet), every race is positively split,
sign-contradicting the reserve-coupled mechanism; the mean observed shape
lies within 0.3 percentage points of the front-loaded model family; and the
model ranking moves with the assumed dive-start credit across its
literature-anchored uncertainty band, making the start correction, rather
than the fatigue physics, the decisive open quantity. The registered
comparison on the expanded dataset, including fitted fatigue parameters and
the held-out evaluation, is reported in
Sections~\ref{sec:results-empirical}--\ref{sec:discussion}.
\end{abstract}
